
\chapter{Что записано на карте}\label{ch:card-data}
Платёжная карта содержит данные с~разными функциями. PAN (Primary Account
Number, основной номер карты) направляет транзакцию к~нужному счёту, CVV (Card
Verification Value, код верификации карты) и~динамическая криптограмма EMV
(см.~гл.~\ref{ch:emv}) подтверждают подлинность карты, часть данных нужна
только в~момент авторизации и~после неё подлежит уничтожению.
Путаница между этими группами ведёт к~нарушению PCI~DSS (Payment Card Industry
Data Security Standard, стандарт безопасности данных индустрии платёжных карт)
и~эксплуатируемым уязвимостям.

\section{Структура PAN}\label{sec:card-data-pan}

Цифры, выбитые или напечатанные на~лицевой стороне карты, образуют PAN, главный
идентификатор карты во~всей платёжной инфраструктуре. PAN встречается повсюду:
в~авторизационном сообщении (DE~2~--- Data Element~2, элемент данных~2,
по~ISO~8583; см.~гл.~\ref{ch:iso8583}), в~клиринговом файле, в~маскированном
виде на~чеке, на~экране мобильного банка, в~базе данных эмитента. Структура
PAN определяет маршрут транзакции и~то, какого эмитента увидит карточная сеть.

PAN имеет переменную длину. Современные спецификации платёжных сетей требуют
поддерживать значения вплоть до~19~цифр~\cite{iso-7812-transition,visa-tadg}.
Подавляющее большинство карт Visa, Mastercard и~Мир используют шестнадцать
цифр; American Express традиционно использует пятнадцать; ряд старых карт
Visa~--- тринадцать. Поле ввода, жёстко фиксирующее длину в~шестнадцать цифр,
отвергает валидные карты Amex и~часть Visa.

\highlight{Первые шесть или восемь цифр PAN~--- IIN (Issuer Identification
Number), идентификатор эмитента и~карточной сети}. Дальше идёт номер счёта,
уникальный для держателя у~конкретного эмитента. Последняя цифра контрольная:
она вычисляется по~алгоритму Луна и~ловит случайную ошибку при вводе или
передаче номера.

\subsection*{IIN и~маршрутизация}

По первым цифрам PAN платёжная инфраструктура определяет эмитента и~карточную
сеть и~маршрутизирует транзакцию по~её правилам. Исторически этот префикс
назывался BIN (Bank Identification Number), и~термин до сих пор сохраняется
в~индустрии. Стандарт ISO~7812 закрепил более точное название IIN, потому что
карты выпускают и~банки, и~небанковские
эмитенты~\cite{iso-7812-transition,ansi-iin}. В~разговорной речи
и~в~документации процессингов по-прежнему используется \enquote{BIN}; оба
термина здесь синонимы.

Самая первая цифра PAN называется MII (Major Industry Identifier) и~указывает
на~крупную отраслевую группу~\cite{iso-7812-1}; остальные цифры префикса уточняют эмитента
и~сеть по~диапазонам из регистра ANSI/ISO и~правил карточной
сети~\cite{iso-7812-2,ansi-iin}. Эта цифра кодирует отрасль, под~которую
блок номеров выделили десятилетия назад; принадлежность к~конкретной сети
и~эмитенту несёт весь IIN. Разработчик работает с~актуальными таблицами BIN/IIN
конкретных сетей и~эмитентов; выводить принадлежность из MII в~рабочей системе
нельзя.

Долгое время IIN состоял из шести цифр, и~миллион комбинаций покрывал нужды
эмитентов. Рост числа финтех-компаний, выпускающих карты, и~развитие
токенизации, где каждому токену нужен свой диапазон, исчерпали свободные
диапазоны BIN. Ревизия ISO~7812 2017~года запустила переход на восьмизначный
IIN, и~сегодня отрасль использует этот формат~\cite{iso-7812-1,iso-7812-transition}.
Шестизначные BIN из обращения не выведены: процессинги и~интеграции обязаны
поддерживать оба формата одновременно ещё много лет.

Процессинг, который продолжает обрезать IIN до шести цифр, рискует неверно
определить эмитента. Два разных эмитента или даже две разные сети могут
совпасть по~первым шести цифрам, но различаться на седьмой и~восьмой. Неверно
определённый эмитент даёт неверный маршрут или отказ без объяснения причины.
Поддержка восьмизначных IIN устраняет дефект маршрутизации; системы определения
сети и~маршрутизации должны учитывать её с~самого начала.

Между IIN и~последней цифрой PAN расположены цифры, идентифицирующие конкретный
счёт держателя у~конкретного эмитента. Эмитент назначает их по~внутренним
правилам: последовательно, случайно или с~учётом продуктовой линейки. Стандарт
способ назначения не регламентирует, поэтому средняя часть PAN внешнему
наблюдателю непрозрачна.

\subsection*{Разновидности PAN}\label{sec:card-data-pan-variants}

Когда на~маршруте появляются сетевые токены и~локальные суррогаты, у~PAN
возникают однокоренные термины с~разной семантикой, и~разные источники называют
один и~тот~же объект по-разному.

\textbf{FPAN} (Funding PAN)~--- реальный номер карты, выпущенный эмитентом
и~используемый как источник для токена; в~правилах ПС~Мир тот~же объект
называется \enquote{Номер Карты (FPAN)}. \textbf{DPAN} (Device PAN)~--- сетевой
токен EMVCo для одного устройства; в~правилах НСПК термин пишется как
\enquote{Токен~(DPAN)} и~остаётся вторым именем \enquote{Токена}, а~не его
подтипом для устройства. \textbf{MPAN} (Merchant PAN)~--- сетевой токен для
одного ТСП; канон карточных сетей~--- MDES for Merchants (M4M) у~Mastercard
и~функциональность для~ТСП в~Visa Token Service плюс Token Management Service
от~Visa Acceptance Solutions у~Visa. Российскому банку с~марта~2022 эти
сервисы недоступны, а~у~НСПК аналога нет: Mir~Pay~SDK app2app реализует
сценарий привязки к~устройству (\textit{provisioning}
и~\textit{In-Application e-commerce} поверх DPAN), а~токена для одного ТСП под
именем MPAN в~правилах ПС~Мир не~определено. \textbf{Pseudo PAN}~---
неформальный термин индустрии для суррогатов vault или результатов
\textit{format-preserving encryption} (FPE); в~официальных документах PCI~SSC
используется \enquote{surrogate value}, в~стандартах сетей и~у~НСПК термин
\enquote{pseudo PAN} не~встречается.

\highlight{Один и~тот~же знак \enquote{DPAN}, \enquote{MPAN}, \enquote{pseudo
  PAN} у~EMVCo, PCI~SSC и~у~российских правил расшифровывается по-разному:
смысл термина задаёт источник, универсального определения нет}. Разбор противоречий,
матрица \enquote{термин~$\times$~источник} и~канон употребления в~этой
книге~--- в~разделах~\ref{sec:token-terminology} и~\ref{sec:token-scope}
главы~\ref{ch:tokenization}.

\subsection*{Алгоритм Луна}

Последняя цифра PAN контрольная и~вычисляется по~алгоритму Луна, который
предложил Ханс Петер Лун (Hans Peter Luhn), инженер International Business
Machines~\cite{luhn-patent}. Алгоритм не требует секретной информации: любой,
кто знает номер карты, проверит контрольную цифру самостоятельно. Поэтому
проверка Луна ловит только случайные опечатки, не подделки.

Алгоритм Луна проверяет PAN так:

\begin{enumerate}
  \item Начиная со второй цифры справа и~двигаясь влево, удвой каждую вторую
    цифру. Если результат удвоения больше~9, вычти~9~--- это эквивалентно
    сложению цифр результата, $14 \to 1 + 4 = 5$.
  \item Сложи все цифры, как удвоенные, так и~неудвоенные.
  \item Если сумма кратна десяти, номер проходит проверку Луна.
  \item Если нет, номер содержит ошибку.
\end{enumerate}

Пример~--- номер \texttt{4561 2612 3456 0892}:
на~рис.~\ref{fig:luhn-check} все четыре шага разложены поразрядно. Сумма
результатов~---~67, не кратна десяти, поэтому этот номер не прошёл бы проверку:
он вымышлен для примера. В~реальном PAN контрольная цифра подбирается эмитентом
так, чтобы итоговая сумма оказалась кратна десяти.

\begin{figure}[H]
  \centering
  \includefiguresvg[width=\linewidth]{assets/figures/ch04-card-data/luhn-check}
  \caption{Проверка Луна для PAN \texttt{4561 2612 3456 0892}.}\label{fig:luhn-check}
\end{figure}

Платёжной форме достаточно проверить введённый номер; вычислять правильную
контрольную цифру не требуется. Компактная проверка на Python:

\codefile{Python}{samples/ch04-luhn/luhn.py}

Алгоритм Луна хорошо отсеивает опечатки и~часть типичных перестановок соседних
цифр, но это \emph{не} криптографическая защита. Он не защищает от
злонамеренного подбора номера, не подтверждает существование карты и~не
заменяет аутентификацию.

\highlight{PAN работает как адрес: по~нему система находит эмитента и~счёт.
Подлинность карты и~право держателя инициировать операцию PAN не доказывает}.
Для этого нужны другие данные; исторически первыми появились данные магнитной
полосы. CVV, CVV2, iCVV, dCVV закрывают каждый свой канал атаки на~карту
и~разбираются в~разделе~\ref{sec:card-data-cvv}.

\section{Магнитная полоса}\label{sec:card-data-magstripe}

На обратной стороне карты находится тёмная магнитная полоса (magnetic stripe,
magstripe), стандартизованная в~1970-х~годах; её минимальная ширина
по~ГОСТ~Р~ИСО/МЭК~7811-6~--- 11,89~мм для~дорожек~1 и~2 и~15,95~мм при~использовании
трёх дорожек~\cite{gost-7811-6}. Полоса по-прежнему присутствует на~большинстве карт мира, хотя
её роль сегодня остаточная: чип и~бесконтактное ядро вытеснили магнитное чтение
из приёма. Требования к~самой записи задаёт ISO/IEC~7811--2, а~структуру
финансовых карт и~данных на~дорожках описывает
ISO/IEC~7813~\cite{iso-7811-2,iso-7813}.

Магнитная полоса содержит три дорожки (tracks) с~разным форматом и~плотностью
записи. В~массовых карточных платежах работают первая и~вторая, третья осталась
факультативной.

\subsection*{Track~1 (IATA)}

Первая дорожка, стандартизованная IATA, единственная из трёх, в~которой
допустимы буквенно-цифровые данные. Для финансовых карт ISO/IEC~7813 задаёт
структуру и~состав полей Track~1 и~Track~2, по~которым инициируется
транзакция~\cite{iso-7813,visa-tadg}. Структура записи на Track~1:

\begin{center}
  \texttt{\%B4561261234560892\^{}IVANOV/IVAN\^{}2612\ldots}
\end{center}

Здесь \texttt{\%} обозначает маркер начала, \texttt{B} обозначает код формата;
для финансовых карт он всегда B. Далее идёт PAN, разделитель \texttt{\^{}}, имя
держателя в~порядке фамилия-имя, ещё один разделитель, срок действия в~формате
две цифры года и~две цифры месяца, сервисный код, данные для верификации CVV
и~дискреционные данные эмитента (поле, в~которое эмитент при персонализации
записывает собственные служебные значения); \texttt{?} обозначает маркер конца.
Track~1~--- единственная дорожка с~именем держателя, и~поэтому именно она
подпадает под~законодательство о~персональных данных.

\subsection*{Track~2 (ABA)}

Вторая дорожка, разработанная Американской банковской ассоциацией (ABA), хранит
только цифровые данные: до~40~символов, из которых используются цифры 0--9
и~несколько служебных символов. Track~2 остаётся основной рабочей дорожкой для
карточных платежей. Когда карта проходит через считыватель или терминал читает
магнитную полосу, в~авторизационном сообщении в~DE~35 отправляются данные
второй дорожки~\cite{visa-tadg}.

Структура Track~2 компактнее:

\begin{center}
  \texttt{;4561261234560892=2612\ldots}
\end{center}

Маркер начала \texttt{;}, далее PAN, разделитель \texttt{=}, срок действия
в~формате две цифры года и~две цифры месяца, затем сервисный код из трёх цифр,
определяющих условия использования карты (требуется ли PIN, разрешены ли
международные транзакции и~так далее), CVV, дискреционные данные эмитента
и~маркер конца \texttt{?}. Имя держателя на Track~2 \emph{отсутствует}: только
цифровой идентификатор.

Каждое поле Track~2 выросло из задач банковских операций 1960-х--70-х~годов.
PAN заменил ручную запись номера на~слип, и~терминал читал цифры с~полосы без
кассира. Срок действия ограничил \enquote{вечно действующие} карты и~упростил
перевыпуск: эмитент автоматически закрывал операции по~истёкшим картам.
Сервисный код закодировал политику эмитента для терминала: три цифры говорят
\enquote{карта только для PIN}, \enquote{международный приём запрещён} или
\enquote{онлайн только}. CVV добавил машинный признак верификации поверх
тиснёных цифр; данные на полосе можно сверить с~вычислением эмитента и~поймать
клон с~пустой полосой.

\practicenote{%
  В~шестнадцатеричном дампе авторизационного сообщения поле DE~35
  (Track~2 Data) совпадает с~содержимым второй дорожки, прочитанным
  с~магнитной полосы, либо с~эквивалентными данными, сгенерированными чипом.
  В~одной строке оказываются PAN, срок действия и~CVV. Полные данные Track~2
  классифицируются стандартом PCI~DSS как чувствительные аутентификационные
  данные (sensitive authentication data, SAD): их \emph{запрещено} хранить
  после авторизации, даже в~зашифрованном виде~\cite{pci-faq-1335}.
}

\subsection*{Track~3: редкое наследие}

Третью дорожку разработала ассоциация Thrift Industry под приложения
с~перезаписью данных прямо на~карте~--- баланс хранился на~полосе, без
обращения к~эмитенту за каждой операцией. В~эпоху дорогих телекоммуникаций идея
казалась разумной. Карточные сети предпочли онлайн-авторизацию через
центральную инфраструктуру: запись баланса на~карту открывала держателю
возможность перемагнитить полосу и~увеличить себе остаток.

В~современных карточных платежах Track~3 не участвует. Сетевые правила
оставляют её как факультативную историческую опцию, но массовая приёмная
инфраструктура строится вокруг Track~1
и~Track~2~\cite{iso-7813,mc-security-rules-2016}.

\subsection*{Сервисный код}

Три цифры сервисного кода на~Track~1 и~Track~2 задают условия приёма карты
и~потому влияют на~маршрутизацию и~проверки безопасности. Каждая позиция
действует независимо. Первая описывает интерчейндж и~наличие чипа, вторая~---
требования эмитента к~авторизации, третья~--- допустимые услуги и~требования
к~PIN. Полный набор значений по~позициям задан в~ISO/IEC~7813:2006,
Table~3~\cite{iso-7813}. На~рис.~\ref{fig:service-code} разобран код
\texttt{201} из дампа MTI~0100 в~гл.~\ref{ch:iso8583}: международная карта
с~чипом, обычная авторизация, без ограничений по~услугам и~без обязательного
PIN.

\begin{figure}[H]
  \centering
  \includefiguresvg[width=\linewidth]{assets/figures/ch04-card-data/service-code}
  \caption{Service code на~Track~2: три независимые позиции; жёлтым подсвечены значения примера \texttt{201}.}\label{fig:service-code}
\end{figure}

В~реализациях платёжных сетей сервисный код входит в~число параметров,
по~которым эмитент различает данные магнитной полосы и~данные чипа. Поэтому
магнитное CVV/CVC1 и~чиповое iCVV не обязаны
совпадать~\cite{visa-vsdc-guide,mc-security-rules-2016}.

\section{Коды верификации}\label{sec:card-data-cvv}

Подлинность карты подтверждают коды верификации: CVV/CVC1 на~магнитной полосе,
CVV2 на~пластике для CNP-операций, iCVV в~чипе. Каждый код вычисляется отдельно
по~своему ключу и~входному набору, поэтому компрометация одного канала не
открывает остальные.

\subsection*{CVV на~магнитной полосе}

CVV у~Visa и~CVC1 у~Mastercard~--- трёхзначное значение, записанное
в~дискреционных данных на~Track~1 и~Track~2 магнитной полосы. Эмитент вычисляет
его из данных карты и~собственных секретных ключей; в~расчёт входят PAN, срок
действия и~сервисный код~\cite{mc-security-rules-2016}. CVV на~карте не
печатается и~глазу не виден. Он доступен только при считывании магнитной
полосы.

Терминал считывает полосу и~отправляет данные Track~2 в~авторизационном
сообщении. Эмитент извлекает CVV из полученных данных, сам пересчитывает
ожидаемое значение по~тем же входным параметрам и~своему ключу и~сравнивает.
Совпадение означает, что полоса не подменена; расхождение говорит, что данные
скопированы или подделаны, и~транзакция отклоняется.

\subsection*{CVV2: код для операций без физического предъявления карты}

Три цифры, напечатанные на~обратной стороне карты (у~American Express это
четыре цифры на~лицевой), образуют CVV2, или CVC2 у~Mastercard. Это код для
операций без физического присутствия карты (\textit{card-not-present}, CNP):
покупок в~интернете, заказов по~телефону. CVV2 вычисляется отдельно
от~магнитного CVV/CVC1 и~намеренно отличается от него, хотя оба значения
привязаны к~одной и~той же карте~\cite{mc-security-rules-2016}.

Если злоумышленник перехватил PAN и~срок действия из чека, утёкшей базы или
фишинга, ТСП с~проверкой CVV2 не примет интернет-покупку без трёх цифр
с~обратной стороны карты. Защита не абсолютна: CVV2 можно подсмотреть,
сфотографировать или выманить, но в~цепочке атаки появляется ещё один фактор,
который придётся компрометировать отдельно.

Магнитное CVV защищает от~подделки полосы: без ключей эмитента злоумышленник,
знающий PAN и~срок действия, не вычислит CVV и~не запишет его на~чистую полосу.
CVV2 защищает канал CNP~--- даже владея Track~2, злоумышленник не узнает три
цифры с~обратной стороны, они вычисляются отдельно и~на~полосе не записаны.

Покупатель оплачивает заказ в~интернете, вводит PAN и~срок действия, но
ошибается в~CVV2. Эмитент видит несовпадение и~возвращает код~N (CVV2 Not
Matched) в~DE~44.\footnote{Точное расположение результата проверки CVV2 в~DE~44
  различается между карточными сетями; у~Visa это subfield~5 \enquote{CVV2 Result
Code}; здесь и~далее ссылка на~DE~44 указывает на~поле в~целом.} Если шлюз или
ТСП поддерживает 3-D~Secure (протокол аутентификации держателя для
CNP-операций, см.~гл.~\ref{ch:3ds}), транзакция переводится в~\textit{challenge}
(этап интерактивной проверки), и~держатель подтверждает личность через ACS
(Access Control Server, сервер аутентификации на~стороне эмитента). После
успешной аутентификации эмитент может одобрить операцию даже с~неверным CVV2:
3DS даёт более сильное подтверждение подлинности, чем три цифры на~пластике,
и~ответственность за мошенничество смещается к~эмитенту по~правилам переноса
ответственности (\textit{liability shift},
см.~раздел~\ref{sec:3ds-liability})~\cite{visa-core-rules}. CVV2 и~3DS работают
как две проверки с~разным весом; логика \enquote{CVV2 не совпал, отказать}
верна при недоступном 3DS.

\subsection*{iCVV: защита от клонирования}

Когда карточные сети начали внедрять чипы EMV и~потребовали отличать чиповые
транзакции от магнитных, появился риск: данные Track~2 из чипа можно было
скопировать на магнитную полосу поддельной карты. Ответом стал iCVV
(integrated circuit card verification value; у Mastercard этот механизм
называется Chip CVC), альтернативное \emph{статическое} значение кода
верификации, прошитое в~чип на этапе персонализации и~используемое для
\textit{chip-initiated transactions}, то~есть транзакций, инициированных чтением
чипа~\cite{visa-vsdc-guide,mc-security-rules-2016}. Оно отличается от значения
магнитной полосы, и~по~нему эмитент отделяет данные чипа от~данных полосы.
От~динамического CVC3, о~котором речь дальше, iCVV/Chip CVC отличается тем, что
не пересчитывается от транзакции к~транзакции.

Скопированные из чипа данные Track~2 Equivalent (эквивалент второй дорожки,
  который чиповое приложение возвращает терминалу для передачи в~авторизационном
сообщении), выданные за магнитную полосу, эмитент отклонит при проверке:
значение рассчитано для чипа, а~не для магнитной полосы. Криптография EMV
остаётся отдельной защитой; iCVV решает узкую задачу~--- не дать
\textit{chip-derived data}, то~есть данным, полученным с~чипа, работать как
данным обычной магнитной полосы.

\subsection*{dCVV и CVC3: отличие от криптограммы EMV}

В~обычной транзакции EMV карта генерирует криптограмму, связанную
с~конкретной суммой, терминалом и~параметрами операции. Криптограмма~---
центральный механизм аутентификации в~EMV~\cite{emvco-book3}; динамические коды
семейства dCVV/CVC3 (dynamic CVV, динамический CVV) применяются точечно.
В~\textit{contactless mag-stripe profile} у~Mastercard, где бесконтактное ядро
эмулирует данные магнитной полосы для приёма устаревшими терминалами,
используется CVC3, меняющийся от транзакции к~транзакции~\cite{mc-contactless-security}.
Сводить любую чиповую проверку к~\enquote{dCVV} ошибочно: для контактного
и~бесконтактного EMV основным доказательством подлинности остаётся
криптограмма, а~CVC3~--- частный механизм в~отдельных профилях сетей.
Криптограмма EMV подробно разбирается в~разделе~\ref{sec:emv-cryptogram}.

\section{Что нельзя хранить после авторизации}\label{sec:card-data-storage}

PCI~DSS жёстко разделяет данные карты на~те, что допустимо сохранить после
авторизации, и~те, что должны исчезнуть сразу. Ошибка здесь может стоить
компании права принимать карты~\cite{pci-dss-4}.

\subsection*{Данные держателя карты}

\highlight{Первая категория, данные держателя карты (cardholder data): PAN, имя
держателя, срок действия и~сервисный код}~\cite{pci-faq-1335}. \highlight{Эти
  данные \emph{можно} хранить после авторизации, но только при соблюдении
требований PCI~DSS}. PAN защищается шифрованием, токенизацией или маскированием
в~зависимости от контекста, доступ к~нему жёстко ограничен и~протоколируется.
PCI~DSS отдельно требует делать PAN нечитаемым при электронном хранении; имя,
срок действия и~сервисный код сами по~себе в~это требование \textit{unreadable}
не входят, но при хранении рядом с~PAN остаются частью cardholder data в~том же
контролируемом периметре~\cite{pci-dss-4,pci-faq-1222}.

\subsection*{Чувствительные данные аутентификации}

Вторая категория, SAD (sensitive authentication data): полные данные магнитной
полосы или их эквивалент на чипе, коды верификации карты и
PIN/PIN-блок~\cite{pci-faq-1335}. Эти данные \textbf{запрещено} хранить после
авторизации в~любом виде: ни в~открытом, ни в~зашифрованном, ни в~хешированном,
ни на бумаге, ни в~логах. Требование 3.3.1 PCI~DSS прямо запрещает сохранять
SAD после авторизации, даже в~зашифрованном виде~\cite{pci-dss-4,pci-faq-1280}.

PAN нужен для возвратов, разрешения диспутов и~отчётности; без него, хотя бы
в~токенизированном виде, ТСП не сможет вести бизнес. SAD нужны только в~момент
авторизации, чтобы подтвердить подлинность карты или держателя. После
авторизации их ценность для бизнеса падает почти до нуля, а~ценность для
злоумышленника максимальна. Получив один PAN, злоумышленник попытается провести
CNP-покупку; у~ТСП и~эмитента остаются CVV2, 3DS и~проверки риска мошенничества.
Получив полные данные Track~2 с~CVV, он изготовит клон карты для физического
терминала. Получив PIN, снимет наличные в~банкомате.

\FloatBarrier
\needspace{4\baselineskip}
\importantnote{%
  Частая инженерная причина утечки SAD~--- журналы. Разработчик
  включает подробное логирование авторизационных сообщений при отладке,
  и~полные данные Track~2 вместе с~CVV оказываются
  в~текстовом файле на~сервере, который команда не считает частью
  платёжного периметра и~не защищает по~требованиям PCI~DSS. PCI~DSS
  прямо требует, чтобы при обычном отображении PAN сводился к~BIN
  и~последним четырём цифрам, если нет документированной деловой
  необходимости показывать больше, а~SAD не попадали в~логи~\cite{pci-faq-1492,pci-dss-4}.
}

\section{Карта как идентификатор и~как инструмент аутентификации}\label{sec:card-data-identity-vs-auth}

Платёжная карта совмещает идентификатор счёта и~средство доказательства, что
операцию инициирует законный держатель; в~инженерной практике эти роли нужно
держать врозь.

\textbf{Карта как идентификатор.} PAN~--- номер, по~которому платёжная система
определяет, к~какому счёту в~каком банке относится транзакция. В~этом смысле
PAN аналогичен номеру банковского счёта или IBAN (International Bank Account
Number, международный номер банковского счёта) и~отвечает на~вопросы
\enquote{куда направить запрос?} и~\enquote{с~какого счёта списать деньги?}.
PAN напечатан на~карте, виден невооружённым глазом, попадает в~чеки
в~маскированном виде, фигурирует в~выписках и~отчётах. Идентификатор открыт
по~своей функции: чтобы направить платёж на~нужный счёт, номер должен быть
известен системе.%

\textbf{Карта как инструмент аутентификации.} CVV, CVV2, данные чипа
с~криптограммой, PIN~--- всё это отвечает на~другой вопрос, можно ли доверять
именно этой операции. Задача~--- подтвердить, что транзакцию инициировал
законный держатель карты или, как минимум, тот, кто физически владеет картой.
В~отличие от идентификации, аутентификация работает только пока секретные
данные остаются секретными: утёк PIN~--- атакующий аутентифицируется вместо
держателя.

Исторически карточная индустрия долго смешивала эти две роли. В~эпоху
импринтеров (механических устройств, переносивших оттиск рельефных цифр карты
на~бумажный слип) 1950--1980-х~годов для совершения платежа достаточно было
предъявить карту: оттиск на~слипе одновременно идентифицировал счёт
и~\enquote{аутентифицировал} держателя~--- физическое присутствие карты
считалось достаточным доказательством. Магнитная полоса добавила CVV, но
проблему не решила, потому что данные полосы статичны и~копируются. Чип
с~динамической криптограммой впервые разделил идентификацию и~аутентификацию
на~уровне технологии. PAN по-прежнему идентифицирует счёт, а~криптограмма,
уникальная для каждой транзакции, доказывает, что чип физически присутствовал
в~момент операции.

Токенизация (см.~гл.~\ref{ch:tokenization}) закрепляет это разделение.
Идентификатор-токен становится бесполезным за пределами конкретного контекста
использования, а~аутентификация выполняется криптографическими средствами,
не зависящими от знания номера карты.

Когда кто-то предлагает \enquote{просто проверить номер карты} для
подтверждения платежа, разделите две сущности: открытый идентификатор
и~аутентификационный фактор, который должен быть секретным, одноразовым
и~привязанным к~контексту.
